\begin{helper}
Fix a terminal coordinate set \(U\), a blue support label \(B\), and a red
label \(J\).  Let \(\mathcal T\) be the finite branch-label type, let
\(\mathcal P\) be the finite transcript family, and write
\(\tau=(P_\tau,A_\tau,Z_\tau)\).  For each branch \(\beta\), let
\(\operatorname{blueTrace}(\beta)\) be its complete terminal blue truth
table, let \(R_{\beta,J}\) be the red support decoded from \(J\), and let
\(O_\beta\) be the terminal open-candidate set after the complete blue trace
has been exposed.

Assume branch labels are functional from the complete terminal blue truth
table.  Assume that every complete terminal assignment compatible with a
realized cell has blue truth table \(\operatorname{blueTrace}(\beta)\), and
that the branch choice and prefix transcript are functions of that complete
assignment.  For every realized cell assume
\[
R_{\beta,J}\subseteq U,\qquad
B\cup R_{\beta,J}\subseteq P_\tau,\qquad
P_\tau\cap A_\tau=\emptyset,
\]
\[
Z_\tau\subseteq A_\tau,\qquad Z_\tau\subseteq O_\beta .
\]
Let
\[
\mathsf{Sibling}(A,\beta,\tau,e)
\]
mean that \(A\subseteq U\) is a complete terminal assignment in the present
sibling at the selected-absence split \(e\in Z_\tau\): it agrees with the
fixed prefix and the complete blue trace of \(\beta\), but has \(e\) present.
Assume precisely that this sibling predicate implies
\[
e\in A
\]
and that \(A\)'s blue truth table is \(\operatorname{blueTrace}(\beta)\).
Assume the fixed-support exhaustion statement: whenever \(A\) is compatible
with a realized fixed-\((B,J)\) cell and an open candidate in that branch is
present in \(A\), that candidate lies in the labelled support
\[
B\cup R_{\beta,J}.
\]

Then the whole present sibling is pruned.  Namely, if \((\beta,\tau)\) is a
realized cell, \(e\in Z_\tau\), and
\(\mathsf{Sibling}(A,\beta,\tau,e)\), then \(A\) is compatible with no
realized fixed-\((B,J)\) branch/transcript cell:
\[
\neg\,\mathsf{Comp}(A,\beta',\tau')
\]
for every realized \((\beta',\tau')\).
\end{helper}
\begin{proof}
Fix \(A,\beta,\tau,e\) with \((\beta,\tau)\) realized,
\(e\in Z_\tau\), and
\(\mathsf{Sibling}(A,\beta,\tau,e)\).  Suppose for contradiction that
\(A\) is also compatible with a realized cell \((\beta',\tau')\).

The sibling condition says that \(A\) has the same complete terminal blue
truth table as \(\beta\).  Compatibility with \((\beta',\tau')\) says that
the same assignment \(A\) has blue truth table
\(\operatorname{blueTrace}(\beta')\).  Since branch labels are functional
from this complete blue truth table, \(\beta=\beta'\).  This is exactly the
blue-trace uniqueness supplied by
\(\noderef{FixedSetHistoryCellRedBranchTranscriptCells}\), and it is the
same functionality used by
\(\noderef{FixedSetHistoryCellRedSupportDeterminedByBlueTrace}\) to keep the
fixed-\(J\) red support from becoming a second independent branch choice.

Now view the alleged compatible cell as a cell with the same branch
\(\beta\).  Because \(e\in Z_\tau\) and
\(Z_\tau\subseteq O_\beta\), the coordinate \(e\) is an open candidate for
this branch.  The sibling condition says \(e\) is present in \(A\).  The
fixed-support exhaustion hypothesis applied to the compatible realized cell
therefore gives
\[
e\in B\cup R_{\beta,J}.
\]
But the original realized cell has
\[
B\cup R_{\beta,J}\subseteq P_\tau,\qquad
Z_\tau\subseteq A_\tau,\qquad
P_\tau\cap A_\tau=\emptyset .
\]
Since \(e\in Z_\tau\), the last two inclusions place \(e\) in
\(A_\tau\).  The support conclusion places \(e\) in \(P_\tau\).  This
contradicts \(P_\tau\cap A_\tau=\emptyset\).

The contradiction was obtained for an arbitrary realized
\((\beta',\tau')\).  Hence every complete assignment in the present sibling
at the selected split is incompatible with every realized fixed-\((B,J)\)
cell.  This is a cylinder-level pruning statement: it excludes the entire
present sibling determined by the fixed prefix, not merely the assignment
obtained by flipping one previously chosen compatible terminal assignment.
\end{proof}
